% Part: incompleteness
% Chapter: incompleteness-provability
% Section: godels-paper

\documentclass[../../../include/open-logic-section]{subfiles}

\begin{document}

\olfileid{inc}{inp}{gop}

\olsection{مقایسه با مقالهٔ اصلی گودل}

ارزش دارد زمانی را به بررسی مقالهٔ سال ۱۹۳۱ گودل اختصاص دهیم. مقدمهٔ
آن اندیشه‌هایی را که همین اکنون بررسی کردیم به‌اجمال ترسیم می‌کند. حتی
اگر مقاله را فقط مروری سریع کنید، به‌آسانی می‌توان دید که در هر مرحله
چه می‌گذرد: گودل ابتدا دستگاه صوری~$P$ را توصیف می‌کند (نحو، اصول
موضوع و قواعد اثبات)؛ سپس توابع و روابط بازگشتی اولیه را تعریف می‌کند؛
پس از آن نشان می‌دهد که~$x B y$ بازگشتی اولیه است و استدلال می‌کند که
توابع و روابط بازگشتی اولیه در~$\Th{P}$ بازنمایی می‌شوند. آنگاه، همانند
آنچه در بالا دیدیم، به اثبات قضیهٔ ناتمامیت می‌پردازد. در بخش~۳ نشان
می‌دهد که می‌توان گزارهٔ اثبات‌ناپذیر را جمله‌ای در زبان حساب گرفت. خاستگاه
لم~$\beta$ همین‌جاست؛ همان لمی که ما نیز برای کار با دنباله‌ها در اثبات
بازنمایی‌پذیری توابع بازگشتی در~$\Th{Q}$ به کار بردیم. گودل تا آنجا
پیش نمی‌رود که مجموعه‌ای کمینه از اصول موضوع کافی را مشخص کند، اما اکنون
می‌دانیم که~$\Th{Q}$ از عهدهٔ این کار برمی‌آید. سرانجام، در بخش~۴
طرحی از اثبات قضیهٔ دوم ناتمامیت ارائه می‌کند.

\end{document}
